\section{Truncated Taylor algebra}
Fix $N\ge0$ and let
\begin{equation}
\mathcal T_N=\C[h]/(h^{N+1}).
\label{eq:TN}
\end{equation}

\begin{definition}[Normalized jet]
For an analytic scalar, vector or matrix function $X$,
\begin{equation}
X_r=\frac{1}{r!}X^{(r)}(\theta_0),
\qquad
X(\theta_0+h)=\sum_{r=0}^{N}X_rh^r+O(h^{N+1}).
\label{eq:jetconvention}
\end{equation}
\end{definition}

\begin{lemma}[Closure of the finite q-Racah construction]
\label{lem:closure}
Suppose every denominator has nonzero constant term and each square-root input has nonzero constant term with a fixed local branch.  Addition, multiplication, reciprocal, quotient, square root and finite summation are recursively well-defined in $\mathcal T_N$.  Consequently every entry of \eqref{eq:rawF} has a unique jet of order $N$.
\end{lemma}

\begin{proof}
Addition and multiplication are inherited from the quotient algebra.  If
\[
a(h)=\sum_{r=0}^{N}a_rh^r,
\qquad a_0\ne0,
\]
and $a(h)b(h)=1$, then
\begin{equation}
b_0=a_0^{-1},
\qquad
b_r=-a_0^{-1}\sum_{k=1}^{r}a_kb_{r-k}.
\label{eq:reciprocal}
\end{equation}
If $c(h)^2=a(h)$ and $c_0^2=a_0$ is the selected branch, then
\begin{equation}
c_r=\frac{1}{2c_0}
\left(a_r-\sum_{k=1}^{r-1}c_kc_{r-k}\right),
\qquad r\ge1.
\label{eq:sqrtrecurrence}
\end{equation}
Quotients reduce to multiplication by a reciprocal.  The q-Racah sum is finite, so coefficient extraction commutes with summation.  Repeated application to \eqref{eq:qnumber}--\eqref{eq:rawF} proves the result.
\end{proof}

\begin{remark}[Formal order versus tested order]
Lemma \ref{lem:closure} is an exact arbitrary-finite-order statement conditional on regularity.  The released implementation is cross-validated through matrix derivative order three and logarithmic-speed derivative order two.  The manuscript does not claim all-order numerical validation.
\end{remark}

\section{Generator and logarithmic-speed recurrences}
Write
\begin{equation}
F(\theta_0+h)=\sum_{r\ge0}F_rh^r,
\qquad
F'(\theta_0+h)=\sum_{r\ge0}(r+1)F_{r+1}h^r.
\end{equation}

\begin{theorem}[Recoupling-generator recurrence]
\label{thm:generatorrecurrence}
If
\[
K(\theta_0+h)=\sum_{r\ge0}K_rh^r,
\qquad K=F'F^T,
\]
then
\begin{equation}
K_r=\sum_{a=0}^{r}(a+1)F_{a+1}F_{r-a}^{T}.
\label{eq:Krecurrence}
\end{equation}
If $FF^T=I$ as a formal series, every $K_r$ is skew-symmetric.
\end{theorem}

\begin{proof}
The coefficient identity follows by multiplying the series for $F'$ and $F^T$.  Differentiating $FF^T=I$ gives $F'F^T+FF'^T=0$, hence $K+K^T=0$ and skewness holds coefficientwise.
\end{proof}

Let
\begin{equation}
\omega(h)=\sum_{r\ge0}\omega_rh^r,
\qquad
\omega_r=(K_r)_{21}.
\label{eq:omegajet}
\end{equation}
If $\omega_0\ne0$, put $s=\operatorname{sgn}(\omega_0)$.  On a zero-free simply connected neighborhood, define the holomorphic branch
\begin{equation}
\mathcal L(h)=\Logh(s\omega(h))=\sum_{r\ge0}L_rh^r,
\label{eq:logjet}
\end{equation}
normalized by $\mathcal L(0)=\log|\omega_0|$.

\begin{lemma}[Holomorphic logarithm recurrence]
\label{lem:logrecurrence}
The coefficients satisfy
\begin{equation}
L_0=\log|\omega_0|,
\qquad
L_n=\frac{1}{n\omega_0}
\left(n\omega_n-\sum_{k=1}^{n-1}kL_k\omega_{n-k}\right),
\quad n\ge1.
\label{eq:logrecurrence}
\end{equation}
\end{lemma}

\begin{proof}
From $\mathcal L'=\omega'/\omega$, equivalently $\omega\mathcal L'=\omega'$, compare the coefficient of $h^{n-1}$.
\end{proof}

\begin{corollary}[Real-axis observables]
\label{cor:firstobservables}
On the real axis near the anchor, $\operatorname{Re}\mathcal L=\log|\omega|$, and
\begin{align}
|\omega(\theta_0)|&=|\omega_0|,\\
\partial_\theta\log|\omega|\big|_{\theta_0}&=L_1,\\
\partial_\theta^2\log|\omega|\big|_{\theta_0}&=2L_2.
\end{align}
Equivalently,
\begin{align}
\frac{\omega'}{\omega}&=\frac{K'_{21}}{K_{21}},\\
\partial_\theta^2\log|\omega|
&=\frac{K''_{21}}{K_{21}}-
\left(\frac{K'_{21}}{K_{21}}\right)^2.
\end{align}
\end{corollary}

\section{Low-order q-number checks}
The first three logarithmic derivatives are
\begin{align}
\partial_\theta\log\qnum{n}
&=n\cot(n\theta)-\cot\theta,
\label{eq:d1q}\\
\partial_\theta^2\log\qnum{n}
&=-n^2\csc^2(n\theta)+\csc^2\theta,
\label{eq:d2q}\\
\partial_\theta^3\log\qnum{n}
&=2n^3\csc^2(n\theta)\cot(n\theta)
-2\csc^2\theta\cot\theta.
\label{eq:d3q}
\end{align}
They propagate through q-factorials, triangle factors, the finite q-Racah sum and dimension amplitudes.  These hand-expanded routes provide cross-checks for the unified coefficient engine; they do not constitute a separate q-Racah implementation.
